% !TEX TS-program = pdflatex
% !TeX spellcheck = en_US
% !TEX root = main.tex

The goal of this evaluation is to investigate the applicability of the proposed algorithm to both regular and context-free path querying.
We measured the execution time of the index creation which solves the reachability problem for both kinds of queries.
The execution time for CFPQ was compared with Azimov's algorithm for CFPQ reachability.
We also investigated the practical applicability of the paths extraction algorithm to both regular and context-free path queries.

For evaluation, we used a PC with Ubuntu 18.04 installed.
It has Intel core i7-6700 CPU, 3.4GHz, and DDR4 64Gb RAM.
We only measure the execution time of the algorithms themselves, thus we assume an input graph is loaded into RAM in the form of its adjacency matrix in the sparse format.
Note, that the time needed to load an input graph into the RAM is excluded from the time measurements.

\he{CFPQ Evaluation}

We evaluate the applicability of the proposed algorithm to CFPQ processing over real-world graphs on a number of classic cases and compare them with Azimov's algorithm.
Currently, only a single path version of Azimov's algorithm exists, and we use its implementation using PyGraphBLAS. Note that it is not trivial to compare our results with the state-of-the-art results provided by~\cite{10.1145/3398682.3399163} (Azimov's algorithm) because our algorithm computes significantly more information. While the state-of-the-art solution computes only reachability facts or a single-path semantics, our algorithm computes data necessary to restore all possible paths.

\he{Dataset}

We use CFPQ\_Data\footnote{CFPQ\_Data is a dataset for CFPQ evaluation which contains both synthetic and real-world data and queries \url{https://github.com/JetBrains-Research/CFPQ\_Data}. Access date: 07.07.2020.} dataset for evaluation.
Namely, we use relatively big RDF files and respective same-generation queries $G_1$~(Eq.~\ref{eqn:g_1}) and $G_2$~(Eq.~\ref{eqn:g_2}) which are used in other works for CFPQ evaluation.
We also use the $Geo$~(Eq.~\ref{eqn:geo}) query provided by~\cite{Kuijpers:2019:ESC:3335783.3335791} for \textit{geospecies} RDF.
Note that we use $\overline{x}$ notation in queries to denote the inverse of $x$ relation and the respective edge.
\begin{align}
\begin{split}
\label{eqn:g_1}
S \to & \overline{\textit{subClassOf}} \ \ S \ \textit{subClasOf} \mid \overline{\textit{type}} \ \ S \ \textit{type}\\   & \mid \overline{\textit{subClassOf}} \ \ \textit{subClasOf} \mid \overline{\textit{type}} \ \textit{type}
\end{split}
\end{align}
\begin{align}
\begin{split}
\label{eqn:g_2}
S \to \overline{\textit{subClassOf}} \ \ S \ \textit{subClasOf} \mid \textit{subClassOf}
\end{split}
\end{align}
\begin{align}
\begin{split}
\label{eqn:geo}
S \to & \textit{broaderTransitive} \ \  S \ \overline{\textit{broaderTransitive}} \\
      & \mid \textit{broaderTransitive} \ \  \overline{\textit{broaderTransitive}}
\end{split}
\end{align}
\begin{align}
\begin{split}
\label{eqn:ma}
S & \to \overline{d} \ V \ d \\
V & \to ((S?) \overline{a})^* (S?) (a (S?))^*
\end{split}
\end{align}

Additionally, we evaluate our algorithm on memory aliases analysis problem: a well-known problem which can be reduced to CFPQ~\cite{Zheng:2008:DAA:1328897.1328464}.
To do it, we use some graphs built for different parts of Linux OS kernel (\textit{arch}, \textit{crypto}, \textit{drivers}, \textit{fs}) and the query \textit{MA}~(Eq.~\ref{eqn:ma})~\cite{10.1145/3093336.3037744}.
The detailed data about all the graphs used is presented in Table~\ref{tbl:graphs_for_cfpq}.

{\setlength{\tabcolsep}{0.2em}
\begin{table*}[h]
    \centering
{
\caption{Graphs for CFPQ evaluation: \textit{bt} is broaderTransitive, \textit{sco} is subCalssOf}
\label{tbl:graphs_for_cfpq}
\scriptsize
%\rowcolors{2}{black!2}{black!10}
\begin{tabular}{|l|c|c|c|c|c|c|c|}
\hline
Graph          & \#V       & \#E        & \#sco & \#type &\#bt & \#a  & \#d \\
\hline
\hline
eclass\_514en  & 239 111    & 523 727    & 90 512    & 72 517    &        ---        & ---  & --- \\
enzyme         & 48 815     & 109 695    & 8 163     & 14 989    &        ---        & ---  & --- \\
geospecies     & 450 609    & 2 201 532  & 0         & 89 062    &        20 867     & ---  & --- \\
go             & 272 770    & 534 311    & 90 512    & 58 483    &        ---        & ---  & --- \\
go-hierarchy   & 45 007     & 980 218    & 490 109   & 0         &        ---        & ---  & --- \\
taxonomy       & 5 728 398  & 14 922 125 & 2 112 637 & 2 508 635 &        ---        & ---  & --- \\
\hline
arch           & 3 448 422  & 5 940 484  &      ---     &  ---   &        ---        & 671 295 & 2 298 947 \\
crypto         & 3 464 970  & 5 976 774  &      ---     &  ---   &        ---        & 678 408 & 2 309 979 \\
drivers        & 4 273 803  & 7 415 538  &      ---     &  ---   &        ---        & 858 568 & 2 849 201 \\
fs             & 4 177 416  & 7 218 746  &      ---     &  ---   &        ---        & 824 430 & 2 784 943 \\
\hline
\end{tabular}
}
\end{table*}
}
\he{Results}

We averaged the index creation time over 5 runs for both single-path Azimov's algorithm (\textbf{Mtx}) and the proposed algorithm (\textbf{Tns}) (see Table~\ref{tbl:CFPQ_index}).

{\setlength{\tabcolsep}{0.2em}
  \begin{table*}[h]
    \centering
    \caption{CFPQ evaluation results, time is measured in seconds}
    \label{tbl:CFPQ_index}
   % \rowcolors{4}{black!2}{black!10}
    \small
    \begin{tabular}{| l | c | c | c | c | c | c | c | c |}
      \hline

      \multirow{2}{*}{Name}  & \multicolumn{2}{c|}{$G_1$} & \multicolumn{2}{c|}{$G_2$} & \multicolumn{2}{c|}{\textit{Geo}} & \multicolumn{2}{c|}{\textit{MA}}\\
      \cline{2-9}
                      & Tns    & Mtx    & Tns  & Mtx  & Tns   & Mtx   & Tns     & Mtx \\
      \hline
      \hline
      eclass\_514en   & 0.24   & 0.27   & 0.25 & 0.26 & ---   & ---   & ---     & ---\\
      enzyme          & 0.03   & 0.04   & 0.02 & 0.01 & ---   & ---   & ---     & ---\\
      geospecies      & 0.08   & 0.06   & $<0.01$ & 0.01 & 26.12 & 16.58 & ---     & ---\\
      go-hierarchy    & 0.16   & 1.43   & 0.23 & 0.86 & ---   & ---   & ---     & ---\\
      go              & 1.56   & 1.74   & 1.21 & 1.14 & ---   & ---   & ---     & ---\\
      pathways        & 0.01   & 0.01   & 0.01 & 0.01 & ---   & ---   & ---     & ---\\
      taxonomy        & 4.81   & 2.71   & 3.75 & 1.56 & ---   & ---   & ---     & ---\\
      \hline
      arch            & ---    & ---    & ---  & ---  & ---   & ---   & 262.45  & 195.51  \\
      crypto          & ---    & ---    & ---  & ---  & ---   & ---   & 267.52  & 195.54  \\
      drivers         & ---    & ---    & ---  & ---  & ---   & ---   & 1309.57 & 1050.78 \\
      fs              & ---    & ---    & ---  & ---  & ---   & ---   & 470.49  & 370.73  \\
      \hline
    \end{tabular}
  \end{table*}
}

We can see that while in some cases our solution is comparable or just slightly better than Azimov's algorithm (\textit{enzyme, eclass\_514en, go}), there are cases when our solution is significantly faster (\textit{go-hierarchy}, up to 9 times faster), and when Azimov's algorithm about 1.3 times faster (all memory aliases and \textit{geospecies} with \textit{Geo} query).
Thus we can conclude that our solution is performant enough because Azimov's algorithm is the fastest known practical CFPQ algorithm~\cite{10.1145/3398682.3399163} and the current version, which we use for evaluation, computes index only for single-path semantics, but our algorithm computes much more information (index for all paths) in a comparable time.


Best to our knowledge, the proposed algorithm is the first algorithm that provides information about all paths of interest (Azimov's algorithm computes information about only one path).
The direct comparison with other solutions is impossible, and we just estimate the running time of our algorithm for a small number of cases.
Namely, we extract all paths with length not greater than 20 edges between all pairs of vertices from indices created for graphs \textit{go} and \textit{eclass\_514en} and query $G_1$.
Paths extraction for one pair of vertices requires 2.64 seconds averaged over all pairs for \textit{go} graph. The~maximal time is 4699 seconds and 217 737 paths were extracted during this time. The average number of paths between two vertices is 184.
For \textit{eclass\_514en} paths extraction for one pair of vertices requires 1.27 seconds averaged over all pairs. The~maximal time is 8.04 seconds and only one path is extracted during this time. The average number of paths between two vertices is 3.
We can see that paths can be extracted in a reasonable time, but a detailed analysis of paths extraction algorithm performance depends on graphs structure.

%Comparison of paths extraction is presented in Figures~\ref{fig:geo_matrix_cfpq} and~\ref{fig:geo_tensors_cfpq}.
%While both methods demonstrate linear time on the length of the extracted path, our generic solution is more than 1000 times slower than Azimov's single path extraction procedure.
%We conclude that current generic all-path extraction procedure is not optimal for single path extraction.

\he{Conclusion}

We conclude that the proposed algorithm is applicable to real-world data processing: the algorithm allows one to solve both the reachability problem and to extract paths of interest in a reasonable time.
While index creation time (reachability query evaluation) is comparable with other existing solutions, the paths extraction procedure should be improved in the future. However, the state-of-the-art solution computes only reachability facts or a single-path semantics, whereas our algorithm computes data necessary to restore all possible paths (all-paths semantics).
Finally, a detailed comparison of the proposed solution with other algorithms for CFPQ and RPQ is required.

To summarize the overall evaluation, the proposed algorithm is applicable to both RPQ and CFPQ over real-world graphs.
Thus, the proposed solution is a promising unified algorithm for both RPQ and CFPQ evaluation.
